% !PRINTED 208
\weberpage{208}
If $I$ is a \gesperrt{primitive} function of the $a$, that is, one whose
numerical coefficients have no common divisor, then the coefficients of the
functions $\Phi$ can in any case have no common divisor other than a power
of 3.

What we have to prove is that they are all divisible by $3^\nu$, or, what is
the same thing, that if we assume $\Phi(A, D)$ to be a primitive function,
then $\nu = 0$ must hold.

The question is therefore this: Can functions with the divisor 3 arise from
the primitive function $\Phi(A, D)$ or $B\Phi(A, D)$ when $A$, $B$, $D$ are
replaced by their expressions in the $a$? Since $B$ is primitive, by
\S.\,2 the function $B\Phi(A, D)$, expressed in the $a$, can have the
divisor 3 only if $\Phi(A, D)$ has it.

The question is thus reduced to this: Can the primitive function
$\Phi(A, D)$ acquire the divisor 3 when $A$, $D$ are replaced by their
expressions?

That this question must be answered in the negative we can easily see as
follows. We first imagine all those terms removed from $\Phi(A, D)$ whose
coefficients are divisible by 3; for evidently the remaining terms must
still retain the same property of acquiring the divisor 3 upon substitution
of the expressions for $A$, $D$.

Secondly, we may assume that $\Phi(A, D)$ does not have the factor $D$, for
by omitting this factor the property in question would, by the theorem
already mentioned (\S.\,2), not be destroyed.

But then a number divisible by 3 would have to arise from $\Phi(A, D)$ when
all the $a$, with the exception of $a_2$, are set equal to zero and
$a_2 = 1$. This gives $D = 0$, $A = 1$, and so $\Phi(1, 0)$ would have to be
a number divisible by 3; but this is the coefficient of the highest power of
$A$ in $\Phi(A, D)$, which by hypothesis is not divisible by 3.
With this we have therefore proved:

\gesperrt{Every integral invariant of the biquadratic form can be
represented rationally and with integer coefficients by $A$, $B$, $D$, and
indeed in one of the two forms}
\[
\Phi(A, D), \quad B\Phi(A, D).
\]
